\originalchapter{激光、增益与共振腔}
\originalsection{受激发射与反转}
本章以6.974的激光专题为补充, 解释用于原子控制的光源. 两能级在同一频率的吸收和受激发射系数相同; 仅用一个连续泵浦通常无法把稳态激发布居提高到$1/2$以上. 三能级或四能级激光通过不同跃迁输运布居, 使激光上态相对下态反转.

设反转数$D=N_u-N_l$, 腔光子数$n$, 泵浦建立反转速率$P$, 反转弛豫$\gamma_\parallel$, 腔损耗$\kappa$, 有效受激耦合常数$G$. 简化速率方程为$\dot n=(GD-\kappa)n+\beta\gamma_\parallel N_u$,$\dot D=P-\gamma_\parallel D-2GDn$. 初始自发项提供种子, 但阈值由$GD=\kappa$决定. 阈上稳态反转近似钳制在$D_{\rm th}=\kappa/G$, 多余泵浦主要增加$n$. 系数$2$来自每次受激跃迁既减少上态又增加下态; 若把$D$定义成上态布居而不是差, 方程系数也随之改变.

\originalsection{腔模、纵模选择与线宽}
腔往返相位$2kL+\phi_{
m mirrors}=2\pi m$决定纵模, 理想空腔自由光谱范围$c/(2L)$. 增益谱宽决定可振荡纵模范围; 腔、增益与色散共同决定实际频率. 线性化$n,D$的阈上方程会出现耦合振荡, 当泵浦与损耗在不同时间尺度上响应时产生弛豫振荡. 这与单原子的Rabi振荡有不同的变量和耗散机制.

相干激光的相位仍可因自发发射与技术噪声扩散. 相位做随机游走会使$g^{(1)}(\tau)$指数衰减, 因而产生Lorentz线. 在简单单模理想模型中, 相位扩散线宽随平均光子数增加而降低; 实际频率噪声还包括腔长波动、泵浦噪声和幅相耦合. 精密光谱必须比较原子线宽、激光线宽和观测时间, 不能把理想相干态的Poisson统计直接理解为绝对单色.
